\section{Local-group typings and the semilattice boundary}
\label{app:local-group}

\paragraph{Semigroup convention.}
Throughout this appendix $\mathbf{LG}$ is tested on the
\emph{positive image} $S=\widehat{s}(\Sigma^+)$ of the typing,
not on the canonical completion $S^\sharp$.
Auxiliary group-valued maps and the displayed monoid morphisms
are implementations whose restrictions to nonempty words define
the semantic bias (Lemma~\ref{lem:semigroup-transfer}).

$\mathbf{LG}$ is the pseudovariety of finite semigroups $S$ in which every
local monoid $eSe$, with $e$ idempotent, is a group; classically
$\mathbf{LG}=\mathbf G*\mathbf D$~\cite{Straubing1985,Tilson1987}.  It is
exactly the class of finite semigroups containing no copy of the two-element
semilattice $U_1=\{1,\zeta\}$ ($\zeta^2=\zeta$)
\cite[Proposition~2.3]{AlmeidaEtAl2009}; that semilattice is the positive
image of the typing behind the exponential lower bound of~\cite{kuriyamaCFG}.
Theorem~\ref{thm:lg-characterization} adds to this \emph{known} algebra the
sliding-factor group typing used for the MCFG quantitative bound.

Throughout, $\Sigma$ is nonempty, $k\ge1$, and $A$ is a finite group.  For a
word $w$ with $|w|\ge k+1$ let $\sigma_k(w)$ be the sequence of its factors of
length $k+1$, ordered by their starting positions.

\begin{definition}[Local-group typing]
\label{def:local-group-typing}
Let $g:\Sigma^{k+1}\to A$ be any map.  For $|w|\ge2k$ put
\[
 \gamma_{k,g}(w):=g(F_1)g(F_2)\cdots g(F_{|w|-k}),
 \qquad \sigma_k(w)=(F_1,\ldots,F_{|w|-k}),
\]
the ordered product in $A$, and put $\gamma_{k,g}(w):=1_A$ for
$|w|<2k$.  The \emph{local-group typing} is
\[
 \lambda_{k,g}(w):=\bigl(\widehat{s}^{k,k}_\Sigma(w),\gamma_{k,g}(w)\bigr).
\]
A finite typing $\widehat{s}$ is \emph{local-group-determined} if
$\widehat{s}\preceq\lambda_{k,g}$ for some $k$, $A$ and $g$.
\end{definition}

For $\ell=k$, words of length $<2k$ are \emph{short} and recorded exactly,
while a \emph{long} $w$ is recorded by
$\langle\operatorname{pref}_k(w),\operatorname{suf}_k(w)\rangle$
(Definition~\ref{def:prefix-suffix-typing}).  For $s,p\in\Sigma^k$ write
$J_g(s,p):=\prod_{i=1}^{k}g\bigl((sp)[i..i+k]\bigr)$, the ordered product over
the $k$ factors of length $k+1$ of $sp$ that start in~$s$.

\begin{lemma}[Local-group typings are morphisms]
\label{lem:local-group-morphism}
The kernel of $\lambda_{k,g}$ is a congruence of finite index on
$\Sigma^*$.  Hence $\lambda_{k,g}$ is a homomorphism onto a finite
monoid $M_{k,g}$.  For long $u,v$,
\[
 \lambda_{k,g}(uv)
 =\bigl(\langle\operatorname{pref}_k u,\operatorname{suf}_k v\rangle,\;
   \gamma_{k,g}(u)\,J_g(\operatorname{suf}_k u,\operatorname{pref}_k v)\,
   \gamma_{k,g}(v)\bigr).
\]
\end{lemma}

\begin{proof}
The window component is a morphism by
Definition~\ref{def:prefix-suffix-typing}.  For long $u,v$, the factors of
length $k+1$ of $uv$ are, in order, those of $u$, the $k$ factors starting in
the last $k$ positions of $u$ (inside
$\operatorname{suf}_k(u)\operatorname{pref}_k(v)$), and those of $v$.  If $u$
is short and $v$ long, the factors starting in $u$ lie inside
$u\operatorname{pref}_k(v)$ and the rest are those of $v$, so
$\gamma_{k,g}(uv)$ is determined by $u$, $\operatorname{pref}_k(v)$ and
$\gamma_{k,g}(v)$; the case of long $u$ and short $v$ is symmetric, and if
both are short $uv$ is determined exactly.  Thus $\lambda_{k,g}(uv)$ is a
function of $\lambda_{k,g}(u)$ and $\lambda_{k,g}(v)$, and the image has at
most $|\Sigma|^{<2k}+|\Sigma|^{2k}|A|$ elements.
\end{proof}

\begin{lemma}[Block decomposition]
\label{lem:local-group-blocks}
Let $w=x_0\beta_1x_1\cdots\beta_qx_q$ be a factorization of a word into
designated blocks $\beta_j$ and fixed words $x_j$.  Let
$w'=x_0\beta'_1x_1\cdots\beta'_qx_q$, where
$\widehat{s}^{k,k}_\Sigma(\beta'_j)=\widehat{s}^{k,k}_\Sigma(\beta_j)$ for every $j$.
Then $\widehat{s}^{k,k}_\Sigma(w')=\widehat{s}^{k,k}_\Sigma(w)$.  If $w$ is long, there are
$c_0,\ldots,c_q\in A$, depending only on the $x_j$ and the window types
of the $\beta_j$, such that
\[
 \gamma_{k,g}(w)=c_0\,\gamma_{k,g}(\beta_1)\,c_1\cdots
   \gamma_{k,g}(\beta_q)\,c_q,
 \qquad
 \gamma_{k,g}(w')=c_0\,\gamma_{k,g}(\beta'_1)\,c_1\cdots
   \gamma_{k,g}(\beta'_q)\,c_q.
\]
\end{lemma}

\begin{proof}
A short block is recorded exactly, so $\beta'_j=\beta_j$, and we absorb it into
the fixed words.  Every remaining block is long and its replacement is long
with the same first and last $k$ letters; hence $w'$ is long exactly when $w$
is, with the same first and last $k$ letters, which gives the window claim.  A
factor $F$ of length $k+1$ of $w$ either lies inside a long block $\beta_j$ or
meets each long block it touches only in a prefix or a suffix of length at
most $k$ (meeting both ends would force $|\beta_j|\le k-1$).  Factors of the
second kind are spelled by the fixed words and the windows of the blocks, and
occur in the same order in $w'$.  Ordering all factors by starting position,
those inside $\beta_j$ form one consecutive run contributing
$\gamma_{k,g}(\beta_j)$; let $c_j$ be the ordered product of the remaining
factors between consecutive runs.
\end{proof}

\begin{theorem}[Characterizations of local-group positive images]
\label{thm:lg-characterization}
For a finite typing $s:\Sigma^+\to\mathsf S$ put $S:=\widehat{s}(\Sigma^+)$.
The following are equivalent.
\begin{enumerate}[label=\textup{(\roman*)},leftmargin=*]
\item $S\in\mathbf{LG}$; that is, $eSe$ is a group for every idempotent
$e$ of $S$.
\item $S$ has no subsemigroup isomorphic to $U_1$.
\item Every idempotent of $S$ lies in the minimal ideal of $S$.
\item $\widehat{s}$ is local-group-determined.  One may take $k=|S|$.
\end{enumerate}
\end{theorem}

\begin{proof}
We include elementary arguments for the standard equivalence of (i)--(iii)
(cf.~\cite[Proposition~2.3]{AlmeidaEtAl2009}), independently of the new
sliding-factor statement (iv).  Let $K$ be the minimal ideal of $S$.

(i)$\Rightarrow$(ii).  If $\{e,z\}\cong U_1$ with
$ez=ze=z\ne e$, then $z$ is a nonidentity idempotent in the group
$eSe$, which is impossible.

(ii)$\Rightarrow$(iii).  For an idempotent $e\in S$ choose $a\in K$
and put $z=(eae)^\omega$.  Then $z$ is an idempotent in
$K\cap eSe$ with $ez=ze=z$.  If $z\ne e$, the pair
$\{e,z\}$ is a copy of $U_1$, a contradiction.  Hence $e=z\in K$.

There is also a direct proof of (i)$\Rightarrow$(iii) that does
not invoke the two-element semilattice: for any idempotent $e$ and
$a\in K$, the element $eae$ lies in $K\cap eSe$; since $eSe$
is a group with identity $e$, it has an inverse $y\in eSe$.
Consequently $e=(eae)y\in K$, because $K$ is an ideal.

(iii)$\Rightarrow$(i).  The minimal ideal of a finite semigroup
is completely simple~\cite{Howie1995}.  For an idempotent $e\in K$,
the ideal property gives $eSe=eKe$; this is the maximal subgroup
(the $\mathcal H$-class) containing $e$, so it is a group.

(iii)$\Rightarrow$(iv).  By the Rees theorem~\cite{Howie1995} write
$K=\mathcal M(H;I,\Lambda;P)$, with product
$(i,a,\lambda)(j,b,\mu)=(i,ap_{\lambda j}b,\mu)$, and put $n:=|S|$.
\emph{Long products lie in $K$.}  The Rees quotient $S/K$ has no idempotent
other than $0$.  For $s_1,\ldots,s_m\in S$ with $m:=|S/K|\le n$, if no prefix
product in $S/K$ were $0$, two would coincide, $p_i=p_j$ with $i<j$, and then
$p_i=p_ix^\omega=0$ for $x=s_{i+1}\cdots s_j$.  Hence $\widehat{s}(w)\in K$ whenever
$|w|\ge n$.
\emph{Right action.}  Fix $i_0\in I$ and for $\lambda\in\Lambda$ let
$e_\lambda:=(i_0,p_{\lambda i_0}^{-1},\lambda)$, an idempotent with
$xe_\lambda=x$ for every $x=(i,a,\lambda)\in K$.  Writing
$e_\lambda s=(j,c,\mu)\in K$ for $s\in S$ gives
\[
 (i,a,\lambda)\,s=(i,\;a\,\rho(\lambda,s),\;\mu(\lambda,s)),
 \qquad \rho(\lambda,s):=p_{\lambda j}c,
\]
with $\rho(\lambda,s)\in H$ and $\mu(\lambda,s)\in\Lambda$ depending only on
$(\lambda,s)$.  Right multiplication by $s$ preserves the $I$-coordinate, as is
explicit in this formula.  Dually, for $z=(i,a,\lambda)\in K$
choose an idempotent $f_i\in K$ with $f_i z=z$.
Then $sz=(sf_i)z$ and $sf_i\in K$, so the Rees product shows that
left multiplication preserves the $\Lambda$-coordinate.  In
particular, for $|w|\ge n$ the $\Lambda$-coordinate of $\widehat{s}(w)$ is
that of $\widehat{s}(\operatorname{suf}_nw)$, and the $I$-coordinate that of
$\widehat{s}(\operatorname{pref}_nw)$.
\emph{Window factorization.}  Let $w=a_1\cdots a_m$ with $m\ge2n$, and let
$\lambda_j$ be the $\Lambda$-coordinate of $\widehat{s}(a_{j-n+1}\cdots a_j)$ for
$j\ge n$.  Iterating the right action gives
\[
 \widehat{s}(w)=\Bigl(i,\;a\prod_{j=n}^{m-1}\rho(\lambda_j,\widehat{s}(a_{j+1})),\;
 \lambda_m\Bigr),
 \qquad \widehat{s}(a_1\cdots a_n)=(i,a,\lambda_n),
\]
and the $j$-th factor depends only on the factor $a_{j-n+1}\cdots a_{j+1}$ of
length $n+1$.  With
$g(F):=\rho\bigl(\lambda(\widehat{s}(F[1..n])),\widehat{s}(F[n+1])\bigr)$ for
$F\in\Sigma^{n+1}$, the product is $\gamma_{n,g}(w)$.  So $\widehat{s}(w)$ is determined
by $\operatorname{pref}_nw$, $\operatorname{suf}_nw$ and $\gamma_{n,g}(w)$;
words of length $<2n$ are recorded exactly by $\widehat{s}^{n,n}_\Sigma$, and $\eps$ is
alone in its $\lambda_{n,g}$-class.  Thus $\widehat{s}\preceq\lambda_{n,g}$ with $A=H$.

(iv)$\Rightarrow$(i).  $\mathbf{LG}$ is closed under homomorphic images, so it
suffices to show $T:=\lambda_{k,g}(\Sigma^+)\in\mathbf{LG}$.  A short element
is never idempotent, since the square of a short word is longer.  Let
$e=(\langle p,s\rangle,\gamma_e)$ be idempotent.  By
Lemma~\ref{lem:local-group-morphism}, $\gamma_eJ_g(s,p)\gamma_e=\gamma_e$, and
for every $x\in T$, $exe$ has window $\langle p,s\rangle$.  Hence $eTe$ lies in
$H_e:=\{(\langle p,s\rangle,\alpha):\alpha\in A\}$, whose product
$\alpha*\beta=\alpha J_g(s,p)\beta$ makes $\alpha\mapsto\alpha J_g(s,p)$ an
isomorphism onto $A$.  Thus $eTe$ is a finite subsemigroup of a group, hence a
group.
\end{proof}

\begin{proposition}[Window--group typings are local-group typings, strictly]
\label{prop:li-g-strict}
\begin{enumerate}[label=\textup{(\roman*)},leftmargin=*]
\item If $\widehat{s}\preceq(\widehat{s}^{k,\ell}_\Sigma,g)$ for a homomorphism $g$ into a
finite group, then $\widehat{s}\preceq\lambda_{K,g'}$ with $K=\max\{k,\ell,1\}$
and $g'(F):=g(\text{last letter of }F)$.  Hence
$\mathbf{LI}\vee\mathbf G\subseteq\mathbf{LG}$.
\item Let $\Sigma=\{\mathtt a,\mathtt b,\mathtt c\}$.  Let
$h_{\mathtt{ab}}:=\lambda_{1,g}$ with $g(\mathtt{ab})=1\in\mathbb Z_2$
and $g(F)=0$ otherwise; it records the parity of the number of
occurrences of $\mathtt{ab}$, together with the window.  Then
$h_{\mathtt{ab}}$ is not window--group-determined.  Hence
$h_{\mathtt{ab}}(\Sigma^+)\in\mathbf{LG}\setminus(\mathbf{LI}\vee\mathbf G)$.
\end{enumerate}
\end{proposition}

\begin{proof}
(i) For $|w|\ge2K$,
$g(w)=g(\operatorname{pref}_Kw)\,\gamma_{K,g'}(w)$, and shorter words are
recorded exactly.  For the inclusion, let $S\in\mathbf{LI}\vee\mathbf G$ and
take $\Sigma=S$ with the evaluation typing $s:\Sigma^+\twoheadrightarrow S$.
Lemma~\ref{lem:mcfg-li-group-window} makes $\widehat{s}$ window--group
determined, hence local-group-determined by the first part, so
$S\in\mathbf{LG}$ by Theorem~\ref{thm:lg-characterization}.

(ii) Suppose $h_{\mathtt{ab}}\preceq(\widehat{s}^{k,\ell}_\Sigma,g')$ with
$g':\Sigma^*\to A'$.  Let $K=\max\{k,\ell,1\}$ and let $q$ be an even
multiple of the exponent of $A'$.  Put
$u=\mathtt a^K(\mathtt{abc})^q\mathtt a^K$ and
$v=\mathtt a^K(\mathtt{bac})^q\mathtt a^K$.  Both are long, with prefix
$\mathtt a^k$ and suffix $\mathtt a^\ell$, and
$g'(u)=g'(\mathtt a)^{2K}=g'(v)$.  But $u$ contains $q$ occurrences of
$\mathtt{ab}$, while $v$ contains exactly one, at
$\mathtt a^K\mathtt b$, so $h_{\mathtt{ab}}(u)\ne h_{\mathtt{ab}}(v)$.
Membership in $\mathbf{LG}$ follows from
Theorem~\ref{thm:lg-characterization}; if the positive image were in
$\mathbf{LI}\vee\mathbf G$, Lemma~\ref{lem:mcfg-li-group-window} would make
$h_{\mathtt{ab}}$ window--group-determined.
\end{proof}

\begin{lemma}[Short tuples with prescribed local-group type]
\label{lem:local-group-shortest}
Let $G=(V,\Sigma,P,S,\fan)$ be a reduced good MCFG of fan-out at most
$f$ and rule rank at most $r_{\max}\ge1$, and let $\lambda=\lambda_{k,g}$.
Let $b_{\mathrm w}:=f(2k+1)n_{\mathrm w}(r_{\max}+1)\bar\theta_G$, where
$n_{\mathrm w}\le|V||M_{k,k}|^f$ is the number of non-start typed
nonterminals of the trimmed $\widehat{s}^{k,k}_\Sigma$-refinement.  Put
$n_\lambda:=|V||M_{k,g}|^f$, and let $M_{A,f}$ be the constant of
Section~\ref{subsec:group-invariants}.  For every $B\in V$ and every
realizable $\mathbf t\in M_{k,g}^{\fan(B)}$ there is $\tuple u\in L_B(G)$
with $\lambda^{(\fan(B))}(\tuple u)=\mathbf t$ and
\[
 \|\tuple u\|_1\le\beta_\lambda:=(r_{\max}+1)\,b_{\mathrm w}
 \max\{1,(M_{A,f}-1)\,n_\lambda\}.
\]
If $A$ is abelian, $M_{A,f}-1$ may be replaced by $\mathsf d(A^f)$.
\end{lemma}

\begin{proof}
Let $T$ be minimal as in Lemma~\ref{lem:heavy-node-count} with root
$\lambda$-vector $\mathbf t$ and threshold $b=b_{\mathrm w}\ge\bar\theta_G$.
For a heavy node $v$ labelled $C$, with yield $\tuple x$ and window vector
$\boldsymbol\omega=\widehat{s}^{k,k}_\Sigma(\tuple x)$, the typed nonterminal
$C_{\boldsymbol\omega}$ is retained in the trimmed $\widehat{s}^{k,k}_\Sigma$-refinement,
because $T$ can be inserted into a successful derivation ($G$ is reduced).
Lemma~\ref{lem:boundary-short-canonical} (trivial $\kappa$) gives a
replacement tuple of $C$ with the same window vector and length at most
$b_{\mathrm w}$.  For pairwise incomparable heavy nodes $v_1,\ldots,v_m$ and
$J\subseteq[m]$, let $T_J$ replace the subtrees at $v_i$, $i\in J$.  The
components of the $v_i$ are disjoint contiguous blocks of the root components.
By Lemma~\ref{lem:local-group-blocks}, every $T_J$ has the root window vector of
$T$; a short root component is unchanged, and on each long root component
$\gamma_{k,g}$ is an ordered product $c_0y_1c_1\cdots y_qc_q$ of block values
between constants that do not depend on $J$.  The long root components thus
form a block-replacement family over $A$
(Definition~\ref{def:block-family}).  If $m\ge M_{A,f}$, or $m>\mathsf d(A^f)$
for abelian $A$, Lemma~\ref{lem:block-replacement} gives a nonempty $J$ with
the root $\lambda$-vector of $T$, while $T_J$ is strictly shorter.  This
contradicts minimality, and Lemma~\ref{lem:heavy-node-count} with
$N=M_{k,g}$ gives the bound.
\end{proof}

\paragraph{Detailed proof of Theorem~\ref{thm:local-group-polydata}.}
The theorem is stated in Section~\ref{subsec:algebraic-typed-data};
we verify here the bound using the sliding-factor characterization
and simultaneous ordered group-block replacement.



\begin{proof}[Proof of Theorem~\ref{thm:local-group-polydata}]
By Theorem~\ref{thm:lg-characterization}, $\widehat{s}\preceq\lambda_{k,g}$ for some
$k,A,g$ depending only on $\widehat{s}$.  For a non-start typed nonterminal
$X=B_{\mathbf s}$ and $\tuple u\in L_X$,
Lemma~\ref{lem:local-group-shortest} gives $\tuple u'\in L_B(G)$ with the same
$\lambda$-vector and $\|\tuple u'\|_1\le\beta_\lambda$, and
Lemma~\ref{lem:type-preserving-replacement} gives $\tuple u'\in L_X$.  The role of the two shortening stages is distinct: the window stage
provides replacements of length $b_{\mathrm w}$ without prescribing their
group values, while the block-replacement lemma selects a nonempty family
of such replacements whose \emph{combined} ordered group products are
unchanged.  No per-subtree group-value preservation is assumed.  Since
$b_{\mathrm w}=O(\|G\|\bar\theta_G)$ and $n_\lambda=O(\|G\|)$,
the heavy-node bound gives $\beta_\lambda=O(\|G\|^2\bar\theta_G)$
for fixed $f,r_{\max},s,\Sigma$.
There are $O(\|G\|)$ typed rules, and
Lemma~\ref{lem:witness-length-template} then yields the displayed cubic
bound on the \emph{total} encoded positive-data length.
\end{proof}

By Proposition~\ref{prop:li-g-strict}(i), this theorem contains the window and
window--group theorems as special cases, with the same exponent three in the
group case.

\begin{corollary}[Pseudovariety dichotomy]
\label{cor:pseudovariety-dichotomy}
Let $\mathbf V$ be a pseudovariety of finite semigroups.
\begin{enumerate}[label=\textup{(\roman*)},leftmargin=*]
\item If $\mathbf V\subseteq\mathbf{LG}$, then for every alphabet,
every typing $\widehat{s}$ with $\widehat{s}(\Sigma^+)\in\mathbf V$, and all fixed
$f,r_{\max}$, characteristic data of total length
$O(\|G\|^3\bar\theta_G)$ exist as in
Theorem~\ref{thm:local-group-polydata}.
\item If $\mathbf V\not\subseteq\mathbf{LG}$, then
$U_1\in\mathbf V$.  At fan-out one there is a typing $h_{\mathrm U}$ with
$h_{\mathrm U}(\Sigma^+)=U_1$ and targets with $\|G\|,\theta_G=O(n)$,
every characteristic sample of which for
$\mathcal B^{\le2}_{1,h_{\mathrm U}}$ contains a word of length
$3\cdot2^n-2$.  The same holds for every set-driven operator admitting
characteristic samples throughout $\Cmcf{1}{h_{\mathrm U}}$.
\end{enumerate}
\end{corollary}

\begin{proof}
(i) is Theorem~\ref{thm:local-group-polydata}.  For (ii), choose
$S\in\mathbf V\setminus\mathbf{LG}$; by
Theorem~\ref{thm:lg-characterization} it has a subsemigroup isomorphic to
$U_1$, and pseudovarieties are closed under subsemigroups.  Take the doubling
family $L_n$, the typing $h_{\mathrm U}$ and the long type-$1$ word $w_n^*$
of~\cite[Proposition~6.4]{kuriyamaCFG}, where $L_n$ is shown to be
Clark--Eyraud substitutable, hence $(1,h_{\mathrm U})$-tuple-substitutable.
The type filter $h_{\mathrm U}^{-1}(\zeta)$ is
$(1,h_{\mathrm U})$-tuple-substitutable by the remark preceding
Proposition~\ref{prop:abelian-coset-filter},
so by Proposition~\ref{prop:abelian-coset-filter}(i) so is
$L_n^-:=L_n\cap h_{\mathrm U}^{-1}(\zeta)=L_n\setminus\{w_n^*\}$.  If
$K\subseteq L_n$ omits $w_n^*$, then $K\subseteq L_n^-$ and
Proposition~\ref{prop:reconstruction-soundness} gives
$L(\mathcal B^{\le2}_{1,h_{\mathrm U}}(K))\subseteq L_n^-$; thus every
characteristic sample contains $w_n^*$.  The set-driven statement is the
distinguishing-pair argument of~\cite[Proposition~6.4(iv)]{kuriyamaCFG} for
the pair $L_n,L_n^-$.
\end{proof}

\begin{remark}[Scope of the dichotomy]
\label{rem:dichotomy-scope}
\begin{enumerate}[label=\textup{(\roman*)},leftmargin=*]
\item Part~(i) holds at every fixed fan-out; part~(ii) is established only at
fan-out one.  The doubling family does not lift directly: $L_2$ is not
$(2,h_{\mathrm U})$-tuple-substitutable, since the type-$1$ pairs
$(\mathtt[,\mathtt a\mathtt])$ and $(\mathtt[\mathtt a,\mathtt])$ share the
context $\mathtt{[[aa]}\,\blank_1\,\mathtt a\,\blank_2\,\mathtt]$ but have
different distributions (checked exhaustively for $n\le3$; under the finer
typing recording each component's set of letters, still with semilattice image,
the failure appears at $n=3$).  Whether part~(ii) holds at fan-out $f\ge2$ is
the pseudovariety form of Open Problem~\ref{op:fanout-two}.
\item The dichotomy concerns pseudovarieties.  For an individual typing $\widehat{s}$
with $\widehat{s}(\Sigma^+)\notin\mathbf{LG}$, the class $\Cmcf{1}{s}$ need not contain
a coded copy of the doubling family, and whether every such typing forces
superpolynomial characteristic data remains open.
\end{enumerate}
\end{remark}